\chapter{Implantable Cardioverter-Defibrillator (ICD) Implantation}
\chaptermark{Implantable Cardioverter-Defibrillator (ICD)}

\section*{Introduction \& Clinical Indications}
Implantable cardioverter-defibrillator (ICD) implantation is a critical surgical procedure aimed at preventing sudden cardiac death in patients at high risk for life-threatening arrhythmias. The surgery involves the placement of a sophisticated electronic device that continuously monitors the heart's electrical activity and delivers therapeutic interventions when necessary. The ICD is typically implanted in a subcutaneous or submuscular pocket, most commonly located in the infraclavicular region of the chest. Through specialized leads, the device connects to the heart, allowing it to detect and treat malignant ventricular arrhythmias, such as ventricular tachycardia (VT) and ventricular fibrillation (VF). By delivering antitachycardia pacing (ATP) or high-energy shocks, the ICD restores normal heart rhythm, significantly reducing the risk of sudden cardiac death in susceptible individuals. This procedure is particularly significant in the context of managing patients with structural heart disease, inherited arrhythmia syndromes, or those who have survived previous cardiac events.

\begin{itemize}
  \item ICD placement
  \item ICD insertion
  \item Defibrillator implantation
  \item Automatic internal cardiac defibrillator (AICD) placement
  \item Cardiac defibrillator surgery
  \item Transvenous ICD implantation
  \item CRT-D implantation (for biventricular devices)
\end{itemize}

The surgical principle of ICD implantation centers on the precise and stable placement of a cardiac rhythm management system capable of continuous, real-time surveillance of the patient's cardiac electrical activity. The rationale is to provide immediate, automated electrical therapy upon the detection of life-threatening ventricular tachyarrhythmias, thereby preventing sudden cardiac death. 

The device's generator, housing a battery and sophisticated microprocessors, is connected to one or more leads that are meticulously positioned within the heart chambers. These leads serve a tripartite function: sensing the heart's intrinsic electrical signals, delivering low-energy pacing pulses for bradycardia or antitachycardia pacing (ATP), and delivering high-energy defibrillation shocks. Advanced algorithms embedded within the ICD continuously analyze the rate, morphology, and onset of detected arrhythmias to accurately differentiate between benign and malignant events, minimizing the risk of inappropriate therapies.

Upon detection of a rapid, potentially fatal ventricular arrhythmia, such as sustained ventricular tachycardia (VT) or ventricular fibrillation (VF), the ICD initiates a tiered therapy approach. This typically begins with ATP, a series of rapid, precisely timed electrical impulses designed to interrupt the re-entrant circuit causing the tachycardia. If ATP is unsuccessful or if VF is detected, the device rapidly charges its capacitors and delivers a high-energy electrical shock directly to the heart muscle, aiming to depolarize a critical mass of myocardial cells simultaneously and allow the heart's natural pacemaker to regain control. The overarching technical goals of the procedure are to achieve stable lead placement with optimal sensing and pacing thresholds, ensure reliable shock delivery, and establish long-term device function to protect the patient from sudden cardiac death.

The concept of an implantable defibrillator was pioneered in the late 1960s by Dr. Michel Mirowski, a cardiologist, and Dr. Morton Mower, an electrical engineer, at Sinai Hospital in Baltimore. Their motivation stemmed from the tragic loss of a colleague to sudden cardiac death, inspiring the vision of an automatic device to terminate ventricular fibrillation. Early designs involved bulky generators and epicardial patches, requiring open-chest surgery. The first successful human implantation of an automatic internal cardiac defibrillator (AICD) occurred on February 4, 1980, at Johns Hopkins Hospital, performed by Dr. Levi Watkins Jr. This landmark event marked a profound shift in the management of sudden cardiac death.

Over the subsequent decades, ICD technology underwent rapid and transformative evolution. The introduction of transvenous leads in the mid-1980s, allowing for lead placement through the venous system without open-chest surgery, revolutionized the procedure, making it significantly less invasive and more widely accessible. Further advancements included the development of dual-chamber ICDs, offering more sophisticated pacing capabilities for bradycardia, and later, biventricular ICDs (CRT-D) for cardiac resynchronization therapy in heart failure patients. Landmark clinical trials, such as MADIT (Multicenter Automatic Defibrillator Implantation Trial) in the 1990s and SCD-HeFT (Sudden Cardiac Death in Heart Failure Trial) in the early 2000s, provided robust, evidence-based support for the efficacy of ICDs in both primary and secondary prevention of sudden cardiac death, firmly establishing their role as a cornerstone therapy in contemporary cardiology.

ICD implantation is indicated for both primary and secondary prevention of sudden cardiac death in carefully selected patient populations. The decision is guided by comprehensive risk stratification and established clinical guidelines.

\begin{itemize}
  \item Secondary prevention for patients who have survived a sudden cardiac arrest due to ventricular tachycardia (VT) or ventricular fibrillation (VF) not caused by a reversible condition.
  \item Treatment for patients with documented sustained ventricular tachycardia (VT) who have structural heart disease and are at high risk for sudden cardiac death.
  \item Primary prevention for patients with severe left ventricular dysfunction, typically a left ventricular ejection fraction (LVEF) of 35\% or less, due to ischemic or non-ischemic cardiomyopathy, who are at high risk for sudden cardiac death despite optimal medical therapy.
  \item Management of patients with certain inherited cardiac conditions, such as long QT syndrome, Brugada syndrome, hypertrophic cardiomyopathy, or arrhythmogenic right ventricular cardiomyopathy (ARVC), who are at high risk for life-threatening arrhythmias.
  \item Provision of cardiac resynchronization therapy with defibrillator function (CRT-D) for patients with symptomatic heart failure, reduced LVEF (typically $\leq$ 35\%), and a wide QRS complex (e.g., $\geq$ 150 ms) who also meet criteria for ICD implantation.
  \item Prevention of sudden cardiac death in patients with unexplained syncope and inducible sustained ventricular arrhythmias during electrophysiology studies, particularly in the presence of structural heart disease.
  \item Patients with non-ischemic dilated cardiomyopathy and LVEF $\leq$ 35\% who are at least 9 months post-diagnosis and on optimal medical therapy.
\end{itemize}

ICD implantation serves as both a primary and secondary therapeutic intervention, depending on the patient's clinical presentation and risk profile. It is considered a primary prevention strategy when implanted in individuals who have not yet experienced a life-threatening ventricular arrhythmia but are deemed to be at high risk for sudden cardiac death based on their underlying cardiac condition. Examples include patients with severe left ventricular dysfunction (e.g., LVEF $\leq$ 35\%) due to ischemic or non-ischemic cardiomyopathy, or those with specific genetic syndromes predisposing them to arrhythmias. In this context, the ICD acts proactively to prevent a first sudden cardiac event, significantly improving long-term survival.

Conversely, ICD implantation is a secondary prevention strategy when it is performed in patients who have already survived a sudden cardiac arrest or experienced documented sustained ventricular tachycardia or fibrillation. For these individuals, the device is implanted to prevent recurrent, potentially fatal arrhythmic events, which carry a high risk of mortality. In both scenarios, the ICD is a cornerstone therapy, often used in conjunction with optimal medical management, lifestyle modifications, and sometimes cardiac rehabilitation. It is rarely a salvage procedure, as its efficacy is maximized when implanted before irreversible cardiac damage or death occurs. The decision to implant an ICD is always guided by comprehensive risk stratification and adherence to established clinical guidelines from major cardiology societies.

\section*{Relevant Surgical Anatomy}
A comprehensive understanding of the relevant thoracic and cardiac anatomy is essential for successful ICD implantation. The device generator is typically placed in a pocket created in the infraclavicular region, usually on the left side, overlying the pectoralis major muscle. This location is chosen for its accessibility, favorable cosmetic outcome, and protection of the device. 

Key venous access points include the cephalic vein, which lies in the deltopectoral groove, and the subclavian vein, which passes beneath the clavicle and anterior to the first rib. The axillary vein, a continuation of the subclavian, can also be utilized. These veins converge to form the brachiocephalic veins, which then unite to form the superior vena cava (SVC), the primary conduit for leads entering the heart.

Within the heart, the leads are advanced into specific chambers. The right atrial (RA) lead is positioned in the right atrium, often with its tip secured in the right atrial appendage or along the interatrial septum. The right ventricular (RV) lead is advanced through the tricuspid valve into the right ventricle, with its tip typically secured in the RV apex or septum. For biventricular ICDs (CRT-D), an additional lead is placed in a coronary sinus branch to pace the left ventricle. The coronary sinus ostium is located in the right atrium, posterior to the tricuspid valve. 

Careful attention must be paid to avoid injury to the phrenic nerve, which runs along the superior vena cava and right atrium, and the brachial plexus, which lies in close proximity to the subclavian vein. The pericardium and myocardium are critical structures, as lead manipulation and fixation occur within these tissues, requiring precision to prevent perforation or tamponade.

\section*{Preoperative Workup \& Preparation}
Thorough pre-procedural preparation is paramount to ensure patient safety and optimize outcomes for ICD implantation. Patients undergo a comprehensive medical evaluation, including a detailed history and physical examination, with a focus on cardiac and pulmonary status.

Specific preparatory steps include:

\begin{itemize}
  \item \textbf{Laboratory Tests:} Standard laboratory tests are ordered, including a complete blood count (CBC), electrolyte panel (especially potassium and magnesium), renal function tests (BUN, creatinine), liver function tests, and coagulation studies (PT/INR, PTT) to assess baseline status and identify any abnormalities.
  
  \item \textbf{Imaging Studies:} A 12-lead electrocardiogram (ECG) is routinely performed. A chest X-ray is obtained to assess lung fields and cardiac silhouette. An echocardiogram is essential to evaluate cardiac structure and function, particularly left ventricular ejection fraction (LVEF).
  
  \item \textbf{Medication Management:} All medications are reviewed. Anticoagulants (e.g., warfarin, direct oral anticoagulants) are typically held for a specified period before the procedure, often with bridging therapy if indicated, under strict medical supervision to minimize bleeding risk. Antiplatelet agents may also be adjusted or held depending on individual patient risk.
  
  \item \textbf{NPO Status:} Patients are instructed to fast (NPO - nil per os) for at least 6-8 hours prior to the procedure to minimize the risk of aspiration during sedation.
  
  \item \textbf{Informed Consent:} Detailed informed consent is obtained, explaining the procedure, its benefits, potential risks, alternatives, and the implications of living with an ICD.
  
  \item \textbf{Cardiac and Anesthesia Clearance:} Patients with significant comorbidities (e.g., severe valvular disease, advanced heart failure, chronic lung disease) may require pre-operative cardiac and anesthesia consultations to assess their fitness for the procedure and optimize their medical status.
  
  \item \textbf{Antibiotic Prophylaxis:} Intravenous antibiotics (e.g., cefazolin) are administered within 60 minutes prior to the skin incision to reduce the risk of device infection.
  
  \item \textbf{Skin Preparation:} The surgical site (chest and shoulder region) is meticulously cleaned and sterilized with an antiseptic solution (e.g., chlorhexidine-alcohol) immediately prior to the procedure. Hair removal, if necessary, is performed with clippers rather than shaving to reduce skin microtrauma.
\end{itemize}

\textbf{Situations requiring correction, shared decision-making, or an alternative plan:}

\begin{itemize}
  \item Active systemic infection or local infection at the proposed implant site, as this significantly increases the risk of device infection, endocarditis, and sepsis, often necessitating device removal.
  
  \item Uncorrectable severe coagulopathy or bleeding disorder that poses an unacceptable risk of hemorrhage during or after the procedure.
  
  \item Short life expectancy (generally less than 1 year) due to non-cardiac causes (e.g., advanced malignancy, severe end-stage organ failure), where the potential benefits of the ICD are unlikely to be realized.
  
  \item Persistent, untreated psychiatric illness or severe cognitive impairment that would prevent the patient from understanding the procedure, complying with follow-up, or managing potential device-related issues, unless a reliable caregiver is available.
  
  \item Patient refusal after thorough discussion of risks and benefits, demonstrating informed decision-making.
  
  \item Reversible causes of ventricular arrhythmias (e.g., acute myocardial ischemia, severe electrolyte imbalance, drug toxicity, hyperthyroidism) that can be treated without the need for an ICD.
\end{itemize}

\textbf{Relative Contraindications \& Precautions:}

\begin{itemize}
  \item Severe tricuspid regurgitation or prosthetic tricuspid valve, which can complicate transvenous lead placement and increase the risk of lead entanglement or damage to the valve.
  
  \item Chronic kidney disease, especially stage 4 or 5, due to potential risks associated with contrast dye used during venography (if performed) and increased risk of infection.
  
  \item Severe pulmonary hypertension, which can increase the risk of right heart injury during lead manipulation and may complicate lead positioning.
  
  \item Significant venous obstruction (e.g., subclavian vein stenosis, superior vena cava syndrome) or congenital venous anomalies, making transvenous access challenging or impossible, potentially requiring alternative access (e.g., epicardial leads).
  
  \item Pregnancy, where radiation exposure from fluoroscopy poses risks to the fetus; careful consideration, lead shielding, and minimization of fluoroscopy time are required if the procedure is deemed essential.
  
  \item Patients with a high burden of inappropriate shocks due to atrial fibrillation or other supraventricular tachycardias, requiring careful device programming, adjunctive antiarrhythmic medications, or ablation procedures.
  
  \item Prior radiation therapy to the chest, which may affect tissue integrity, healing at the implant site, and increase the risk of pocket erosion or infection.
  
  \item Morbid obesity, which can make pocket creation and lead manipulation more challenging, and may increase the risk of infection.
\end{itemize}

\section*{Surgical Technique \& Procedure}
The implantation of an ICD is typically performed in a cardiac catheterization laboratory or a dedicated electrophysiology operating room under strict sterile conditions.

\begin{itemize}
  \item \textbf{Anesthesia and Patient Positioning:} The patient is positioned supine, often with the left arm slightly abducted and externally rotated to expose the infraclavicular region. The procedure is usually performed under local anesthesia (e.g., lidocaine with epinephrine) with conscious sedation (e.g., midazolam and fentanyl) to ensure patient comfort and cooperation. General anesthesia may be utilized for complex cases, pediatric patients, or patient preference.
  
  \item \textbf{Incision and Pocket Creation:} A curvilinear incision, typically 4-7 cm in length, is made in the infraclavicular region, usually 1-2 cm below the clavicle and medial to the deltopectoral groove, most commonly on the left side. A subcutaneous or submuscular pocket is meticulously dissected using blunt and sharp dissection to accommodate the ICD generator. Meticulous hemostasis is achieved using electrocautery.
  
  \item \textbf{Venous Access:} Venous access is obtained, most commonly via the cephalic vein (identified and cannulated through a cutdown in the deltopectoral groove) or the subclavian vein (accessed via percutaneous puncture, often guided by fluoroscopy or ultrasound). The chosen vein provides a conduit for the leads to reach the heart.
  
  \item \textbf{Lead Insertion and Positioning:} Under continuous fluoroscopic guidance, one or more leads are advanced through the vein, superior vena cava, and into the appropriate heart chambers. For a single-chamber ICD, a lead is placed in the right ventricle (RV apex or septum). For a dual-chamber ICD, an additional lead is placed in the right atrium (RA appendage or septum). For a biventricular ICD (CRT-D), a third lead is typically advanced into the coronary sinus ostium and then into a suitable coronary sinus branch to pace the left ventricle. The lead tips are secured, often by active fixation (a helix screwed into the myocardium) or passive fixation (tines lodging in trabeculae).
  
  \item \textbf{Lead Testing:} Once positioned, the leads are thoroughly tested to ensure optimal electrical parameters. This involves measuring sensing (the ability to detect the heart's intrinsic electrical signals, typically $\geq$ 5 mV for RV, $\geq$ 1.5 mV for RA), pacing thresholds (the minimum energy required to consistently stimulate the heart, typically $\leq$ 1.0 V at 0.5 ms pulse width), and lead impedance (resistance to current flow, typically 400-1200 ohms). These measurements confirm proper lead function, stability, and integrity.
  
  \item \textbf{Connection to Generator and Programming:} The leads are then connected to the ICD generator. The device is initially programmed with appropriate detection zones for VT/VF, therapy settings (ATP, shock energy), and bradycardia pacing parameters based on the patient's clinical profile.
  
  \item \textbf{Defibrillation Threshold (DFT) Testing (Optional):} In some cases, a defibrillation threshold test may be performed. This involves inducing a ventricular arrhythmia (typically VF) and delivering a shock to ensure the device can successfully terminate it. This step is increasingly controversial and often omitted in modern practice due to associated risks, potential myocardial stunning, and lack of clear survival benefit in many patient cohorts.
  
  \item \textbf{Generator Placement and Wound Closure:} The ICD generator is placed securely within the created pocket, ensuring it lies flat and comfortably. The incision is then closed in layers using absorbable sutures for subcutaneous and dermal layers, and the skin is closed with sutures, staples, or surgical adhesive. A sterile dressing is applied.
\end{itemize}

No drains are typically left in place unless there is significant bleeding or a large hematoma requiring evacuation. The entire procedure usually takes 1 to 2 hours.

\section*{Complications \& Risks}
While ICD implantation is generally safe, as with any surgical procedure, there are potential risks and complications. These can be broadly categorized as general surgical risks, anesthesia-related risks, and procedure-specific device-related risks.

\textbf{General Surgical Risks:}

\begin{itemize}
  \item \textbf{Bleeding and Hematoma:} Accumulation of blood under the skin at the incision site, which may cause pain, swelling, and rarely require surgical drainage or blood transfusion.
  
  \item \textbf{Infection:} Localized pocket infection, which can be superficial or deep, potentially leading to systemic infection (sepsis) or endocarditis. Device infection often necessitates complete device and lead removal, a complex and risky procedure.
  
  \item \textbf{Pain:} Discomfort at the incision site and around the device pocket, usually manageable with oral analgesics.
  
  \item \textbf{Scarring:} Formation of a visible scar at the incision site, which can vary in appearance.
  
  \item \textbf{Allergic Reaction:} To medications, antiseptic solutions, or device components.
\end{itemize}

\textbf{Anesthesia-Related Risks:}

\begin{itemize}
  \item Allergic reactions to sedatives or local anesthetics.
  
  \item Respiratory depression or aspiration during conscious sedation.
  
  \item Cardiovascular events such as hypotension, bradycardia, or arrhythmias.
  
  \item Nausea and vomiting.
\end{itemize}

\textbf{Procedure-Specific and Device-Related Risks:}

\begin{itemize}
  \item \textbf{Pneumothorax or Hemothorax:} Puncture of the lung (pneumothorax) or a blood vessel (hemothorax) during venous access, leading to air or blood accumulation in the chest cavity, potentially requiring a chest tube insertion.
  
  \item \textbf{Cardiac Perforation or Tamponade:} A rare but serious complication where a lead punctures the heart wall, leading to bleeding into the pericardial sac and compression of the heart (cardiac tamponade), requiring urgent pericardiocentesis or surgical repair.
  
  \item \textbf{Lead Dislodgement or Fracture:} The lead moving out of its intended position or breaking, leading to loss of sensing or pacing, and requiring repositioning or replacement.
  
  \item \textbf{Lead Insulation Failure:} Damage to the lead's outer insulation, leading to electrical malfunction, altered impedance, and potential inappropriate therapies.
  
  \item \textbf{Inappropriate Shocks:} The device delivering a shock when no life-threatening arrhythmia is present, often due to sensing of supraventricular arrhythmias (e.g., atrial fibrillation, sinus tachycardia), T-wave oversensing, or lead malfunction. These are distressing and can impact quality of life.
  
  \item \textbf{Device Malfunction:} Rare instances of hardware or software failure requiring device replacement.
  
  \item \textbf{Venous Thrombosis:} Formation of a blood clot in the vein used for lead insertion, potentially leading to arm swelling, pain, or, rarely, pulmonary embolism.
  
  \item \textbf{Nerve Injury:} Damage to nerves in the shoulder or arm (e.g., brachial plexus, intercostobrachial nerve) during pocket creation or venous access, leading to pain, numbness, or weakness.
  
  \item \textbf{Erosion of the Device:} The generator eroding through the skin, particularly in very thin patients, those with infection, or following trauma to the implant site.
  
  \item \textbf{Radiation Exposure:} Minimal but cumulative exposure from fluoroscopy, though efforts are made to minimize this.
\end{itemize}

\section*{Postoperative Care \& Recovery}
Immediately after ICD implantation, the patient is transferred to a post-anesthesia care unit (PACU) or a telemetry unit for close monitoring. Vital signs, cardiac rhythm, and the surgical incision site are continuously observed for signs of bleeding or hematoma. A chest X-ray is routinely performed within a few hours to confirm lead position and rule out acute complications such as pneumothorax or hemothorax. Pain management is initiated with oral or intravenous analgesics. Patients are typically encouraged to ambulate within 24 hours, with assistance, to prevent complications like deep vein thrombosis. An initial device interrogation is performed by a cardiac electrophysiology nurse or technician to verify lead parameters, battery status, and device programming.

For the first 2-4 weeks post-procedure, patients are advised to limit strenuous activity involving the arm on the implanted side. This includes avoiding heavy lifting (typically over 5-10 pounds), reaching overhead, or excessive arm movement (e.g., golf, swimming) to prevent lead dislodgement or fracture. An arm sling may be recommended for a few days to remind the patient of these restrictions. Wound care instructions are provided, emphasizing keeping the incision clean and dry, and avoiding submerging it in water (baths, swimming pools) until fully healed. Most patients are discharged within 24-48 hours, provided there are no immediate complications and they are hemodynamically stable. Long-term, patients will require regular follow-up device checks, typically every 3-6 months, to monitor device function, battery life, lead integrity, and review any stored arrhythmia events or delivered therapies. Full recovery and return to normal activities, excluding contact sports, usually takes 4-6 weeks.

During ICD implantation, the surgeon's primary findings relate to the successful and stable placement of the device components. These intraoperative findings are meticulously documented in the operative report. Key findings include the chosen venous access site (e.g., cephalic vein cutdown, subclavian vein puncture), the ease or difficulty of lead advancement, the final anatomical position of each lead (e.g., right ventricular apex, right atrial appendage, coronary sinus branch), and the stability of lead fixation (active or passive). 

Crucially, the electrical parameters measured during lead testing---sensing amplitudes, pacing thresholds, and lead impedances---are recorded, confirming optimal lead function and myocardial contact. The integrity of the created subcutaneous or submuscular pocket, ensuring adequate space for the generator and absence of excessive bleeding, is also a significant finding.

Unlike many surgical procedures, routine ICD implantation does not involve the excision of tissue for pathological examination. Therefore, there is typically no pathology report generated from the procedure itself. In rare instances, if a lead is extracted due to infection or malfunction, the lead may be sent for microbiological culture to identify causative organisms. Similarly, if the device pocket becomes infected and requires debridement, tissue samples may be sent for culture. However, for an uncomplicated primary implantation, the "findings" are entirely functional and anatomical, documented through fluoroscopic images and electrical measurements rather than histological analysis.

A normal, successful postoperative course following ICD implantation is characterized by a relatively swift recovery and the absence of significant complications. Patients can expect mild to moderate pain at the incision site for the first few days, which is typically well-controlled with oral analgesics. Swelling and bruising around the device pocket are common and usually resolve within 1-2 weeks. The incision site should heal cleanly, without redness, excessive warmth, or discharge, indicating an absence of infection. Patients are generally able to resume light activities within a few days and gradually increase their activity level over 4-6 weeks, adhering to arm movement restrictions to prevent lead dislodgement.

The device itself should function silently and imperceptibly in the background, continuously monitoring the heart rhythm. Patients should not feel the device unless it delivers therapy. If an appropriate shock is delivered for a life-threatening arrhythmia, the patient will experience a sudden, forceful sensation, often described as a "kick to the chest." In the absence of such events, the expected course is one of improved peace of mind regarding sudden cardiac death, with the patient returning to their baseline functional status. Regular device checks will confirm optimal lead parameters and battery longevity, ensuring the ICD continues to provide effective protection.

Structured postoperative care and long-term follow-up are essential for patients with an ICD to ensure optimal device function and patient well-being.

\begin{itemize}
  \item \textbf{Initial Postoperative Visit:} Typically scheduled 1-2 weeks after implantation for a wound check, removal of sutures or staples (if non-absorbable), and a comprehensive device interrogation to confirm stable lead parameters and initial programming.
  
  \item \textbf{Routine Device Interrogations:} Scheduled every 3-6 months, either in-person at a cardiology clinic or remotely via home monitoring systems. These checks assess battery life, lead integrity, device settings, and review any stored arrhythmia events or delivered therapies.
  
  \item \textbf{Cardiology Clinic Visits:} Regular visits with the cardiologist or electrophysiologist to manage underlying heart conditions, optimize medical therapy, and address any symptoms or concerns related to the device or cardiac health.
  
  \item \textbf{Imaging/Labs:} A chest X-ray may be performed if there are concerns about lead integrity or device position. Blood tests may be ordered as part of routine cardiac care or if specific issues arise (e.g., electrolyte imbalances).
  
  \item \textbf{Activity Resumption and Rehabilitation:} Patients are gradually cleared to resume normal activities, with specific instructions to avoid contact sports or activities that could damage the device or leads. Formal cardiac rehabilitation is often recommended for patients with underlying heart failure or post-myocardial infarction.
  
  \item \textbf{Patient Education:} Ongoing education regarding warning signs to report immediately, such as fever, redness, swelling, or discharge at the incision site, persistent pain, dizziness, syncope, or experiencing a shock. Patients are also advised on electromagnetic interference and how to safely use electronic devices.
\end{itemize}

These comprehensive follow-up measures ensure the ICD continues to function optimally, providing ongoing protection against sudden cardiac death, and addressing any potential issues promptly.

\section*{Outcomes \& Prognosis}
Implantable cardioverter-defibrillator (ICD) implantation has become a widespread procedure globally, with hundreds of thousands of devices implanted annually. The incidence of ICD implantation has steadily increased since its introduction, driven by expanding indications based on robust clinical trial evidence and an aging population with a higher prevalence of cardiovascular disease. The patient population receiving ICDs typically consists of older adults, with a mean age often in the mid-60s, and a slight male predominance, reflecting the higher prevalence of ischemic heart disease and cardiomyopathy in men. 

The most common indications for ICD implantation are primary prevention in patients with severe left ventricular dysfunction (LVEF $\leq$ 35\%) due to ischemic or non-ischemic cardiomyopathy, and secondary prevention in survivors of sudden cardiac arrest or those with documented sustained ventricular arrhythmias. While the procedure-related mortality is exceedingly low, generally less than 1\%, morbidity rates due to complications such as lead dislodgement, infection, pneumothorax, or hematoma range from 3\% to 10\%, highlighting the importance of careful patient selection and meticulous surgical technique. Device infections, though infrequent (1-2\%), are particularly serious, often necessitating complete device and lead extraction.

ICDs can terminate eligible ventricular arrhythmias and reduce sudden-death risk in selected primary- and secondary-prevention populations, but they do not treat the underlying cardiomyopathy or prevent all-cause mortality. Benefits, shocks, complications, and quality-of-life effects require individualized counseling.

In primary prevention cohorts, such as those with severe heart failure and reduced ejection fraction, landmark trials like MADIT II and SCD-HeFT have demonstrated a 20-30\% reduction in all-cause mortality compared to conventional medical therapy alone. Functional outcomes are generally good, with most patients experiencing symptom resolution related to their arrhythmias and maintaining a good quality of life, although some may experience anxiety related to the device or potential shocks. Recurrence of ventricular arrhythmias, leading to appropriate ICD therapies (shocks or antitachycardia pacing), is common, occurring in 20-40\% of patients over several years, but these events are successfully managed by the device. 

Inappropriate shocks, while not life-saving, can be distressing and occur in 10-20\% of patients, often due to sensing of supraventricular arrhythmias or lead issues, and require careful device reprogramming or adjunctive therapies. Prognostic factors influencing long-term outcomes include the severity of underlying heart disease, left ventricular ejection fraction, presence of comorbidities (e.g., renal failure, diabetes), and the frequency of appropriate or inappropriate device therapies. Device longevity typically ranges from 5 to 10 years, after which a generator replacement procedure is necessary.

\section*{References}
\begin{enumerate}
  \item MedlinePlus. Implantable Cardioverter-Defibrillator (ICD) Implantation. U.S. National Library of Medicine; 2024. Available from: \url{https://medlineplus.gov/ency/article/007369.htm}
  
  \item Zipes DP, Libby P, Bonow RO, Mann DL, Tomaselli GF, Braunwald E. Braunwald's Heart Disease: A Textbook of Cardiovascular Medicine. 12th ed. Elsevier; 2022.
  
  \item American College of Cardiology/American Heart Association/Heart Rhythm Society (ACC/AHA/HRS) Guidelines for the Management of Patients With Ventricular Arrhythmias and the Prevention of Sudden Cardiac Death. Circulation. 2018;138(10):e272-e391.
  
  \item Boriani G, et al. ESC Guidelines for the diagnosis and management of ventricular arrhythmias and the prevention of sudden cardiac death. European Heart Journal. 2022;43(40):3997-4126.
\end{enumerate}

\clearpage
